\section{System F confluence and checking: proofs}
\label{sec:annex-p07-system-f-confluence-and-checking}

\begin{annexitem}{annex:p7.system-f.parallel-closure}
 {thm:p07-system-f-parallel-closure}{Reflexivity and both inclusions}
\begin{proof}
Structural induction gives parallel reflexivity: a variable uses clause~1;
term application uses the two inductive reflexivities and clause~2; lambda
abstraction, type abstraction and type application use clauses~3, 4 and~5.

Each ordinary root term-beta contraction is an instance of clause~6 using
reflexivity in its body and argument. Each root type-beta contraction is
an instance of clause~7 using reflexivity in its body. A contraction at
an interior position is propagated by the corresponding constructor clause;
in a term application, use reflexivity in the unaffected child. This
proves \(\to_{\mathsf F}\subseteq\Rrightarrow_{\mathsf F}\).

Conversely, induct on a parallel derivation. Clause~1 is zero ordinary
steps. In clause~2, sequentially perform the finitely many steps supplied
by induction in the function and argument positions. Clauses~3, 4 and~5
place the premise's finite sequence under the same constructor. In
clause~6, first reduce the body and argument to their inductively supplied
parallel reducts, under the enclosing lambda and application; then contract
the displayed root term-beta redex. Clause~7 similarly reduces the body
under its type abstraction and then contracts the root type-beta redex.
These are finite sequences even when the raw term also admits infinite
reductions. Thus parallel reduction is included in ordinary finite
reduction. Taking reflexive transitive closures in the two inclusions
proves equality of those closures.
\end{proof}
\end{annexitem}

\begin{annexitem}{annex:p7.system-f.parallel-substitution}
 {thm:p07-system-f-parallel-substitution}{Every parallel constructor and the mixed substitution cases}
\begin{proof}
The following equalities are the capture-avoiding substitution identities
proved for the same raw syntax in
Annex~\ref{annex:p7.system-f.substitution-reduction}:
\[
 (P[Q/x])\tau=(P\tau)[Q\tau/x],
\]
and, for a type binder \(a\) fresh for \(\tau\)'s domain and images,
\[
 (P[\beta/a])\tau=(P\tau)[\beta\tau/a].
\]
For distinct term variables \(x,z\), with \(z\) fresh for \(N'\),
the term-substitution composition law is
\[
 (P'[Q'/z])[N'/x]
    =(P'[N'/x])[Q'[N'/x]/z].
\]
For \(a\) absent from the free type variables of \(N'\), the mixed
law is
\[
 (P'[\beta/a])[N'/x]=(P'[N'/x])[\beta/a].
\]
These identities include changes in annotations and type arguments. All
bound variables in a finite derivation can be renamed fresh for the
finitely many substitution images under consideration. No free context
variable is renamed.

Prove type-substitution stability by induction on the parallel derivation.
At a variable the substitution leaves the term variable unchanged.
For term application apply induction in both children. For lambda
abstraction apply induction to its body and substitute in the annotation;
the lambda congruence clause then applies with that substituted annotation.
For type abstraction first choose its binder fresh for \(\tau\), apply
induction to the body, and retain the binder. Type application applies
induction to its term and replaces the common type argument by its
substituted type. In clause~6, induction yields \(P\tau\Rrightarrow
_{\mathsf F}P'\tau\) and \(Q\tau\Rrightarrow_{\mathsf F}Q'\tau\).
Parallel root term contraction yields \((P'\tau)[Q'\tau/x]\), equal
to the substituted original contractum by the first identity. In clause~7,
choose the type binder fresh. Induction on the body and parallel root
type contraction yield \((P'\tau)[\beta\tau/a]\), equal to
\((P'[\beta/a])\tau\). This exhausts the seven clauses.

For term substitution, induct on the derivation of
\(M\Rrightarrow_{\mathsf F}M'\), keeping the additional premise
\(N\Rrightarrow_{\mathsf F}N'\) arbitrary. At the variable \(x\)
the desired assertion is exactly that premise; at any other variable use
parallel reflexivity. At a term application, apply induction to both
children and use clause~2. At a lambda abstraction rename its binder
fresh for \(x,N,N'\), apply induction to the body and use clause~3.
At a type abstraction rename its binder fresh for all free type variables
of \(N,N'\), apply induction to the body and use clause~4. At a type
application the type argument contains no term variables, so induction
on its term followed by clause~5 gives the result.

In the root term case write \(M=(\lambda z{:}\alpha.P)Q\) and
\(M'=P'[Q'/z]\). Rename \(z\) away from \(x\) and from the free
term variables of both inserted images. The two premise inductions give
\[
 P[N/x]\Rrightarrow_{\mathsf F}P'[N'/x],\qquad
 Q[N/x]\Rrightarrow_{\mathsf F}Q'[N'/x].
\]
Clause~6 contracts the substituted root to
\((P'[N'/x])[Q'[N'/x]/z]\), which is \(M'[N'/x]\) by the
composition identity. The source is exactly \(M[N/x]\).
This single parallel contraction covers every duplicated occurrence of
the inserted argument, including the case in which the bound variable
does not occur and the argument is erased.

In the root type case write \(M=(\Lambda a.P)[\beta]\) and
\(M'=P'[\beta/a]\), choosing \(a\) fresh for the free type variables
of \(N,N'\). Premise induction gives
\(P[N/x]\Rrightarrow_{\mathsf F}P'[N'/x]\). Clause~7 contracts the
substituted root to \((P'[N'/x])[\beta/a]\). The mixed identity makes
this equal to \((P'[\beta/a])[N'/x]=M'[N'/x]\). There are no other
parallel derivation cases. The proof uses neither typability nor termination.
\end{proof}
\end{annexitem}

\begin{annexitem}{annex:p7.system-f.development-join}
 {thm:p07-system-f-development-join}{All original root and contextual overlap cases}
\begin{proof}
Induct on the size of the original term \(M\), considering all possible
parallel derivations from it. The definition of \(M^\star\) is
structural: substitutions in its output are finite and are not recursively
developed again in that same call.
It is also well defined on alpha classes. Renaming a binder preserves
the original constructor cases, and the recursively developed bodies
and capture-avoiding substitutions respect that renaming. Structural
induction gives the same alpha class for the resulting development.

A variable has only its reflexive parallel derivation, and its development
is itself. If \(M=\lambda x{:}\alpha.P\), the only parallel clause
with that outer source constructor is lambda congruence. Thus
\(N=\lambda x{:}\alpha.P'\), with \(P\Rrightarrow_{\mathsf F}P'\).
Induction gives \(P'\Rrightarrow_{\mathsf F}P^\star\), and lambda
congruence gives the result. A type abstraction has the identical argument
with clause~4. No eta rule can remove either abstraction constructor.

For a term application \(M=PQ\), first suppose \(P\) is not a
lambda abstraction in the original term. The last parallel clause must
be application congruence, giving \(N=P'Q'\) from parallel reducts
of \(P,Q\). Induction and clause~2 give
\(P'Q'\Rrightarrow_{\mathsf F}P^\star Q^\star=M^\star\).
Even if \(P'\) becomes a lambda abstraction, the congruence clause
can retain the newly exposed root redex; contraction of that new redex
is not required by the development definition.

Now suppose \(P=\lambda x{:}\alpha.R\). There are exactly two
possibilities for the last parallel clause from \(M\). A root
contraction gives \(N=R'[Q'/x]\), with \(R\Rrightarrow_{\mathsf F}R'\)
and \(Q\Rrightarrow_{\mathsf F}Q'\). The smaller-term induction
hypotheses and parallel substitution give
\[
 R'[Q'/x]\Rrightarrow_{\mathsf F}R^\star[Q^\star/x]=M^\star.
\]
Alternatively, application congruence retains the root. Because a
parallel reduct of a lambda abstraction is again a lambda abstraction
with the same annotation, it has the form
\(N=(\lambda x{:}\alpha.R')Q'\), with the same kind of premise
reductions. Induction and clause~6 now give its parallel root contraction
to \(R^\star[Q^\star/x]\). Thus contracting and retaining the
original term redex meet at the same development.

For a type application \(M=P[\alpha]\), if \(P\) is not originally
a type abstraction, only type-application congruence applies. Write
\(N=P'[\alpha]\); induction gives
\(N\Rrightarrow_{\mathsf F}P^\star[\alpha]=M^\star\), whether
or not its head becomes a type abstraction later.
If \(P=\Lambda a.R\), a root contraction gives
\(N=R'[\alpha/a]\), where \(R\Rrightarrow_{\mathsf F}R'\).
Induction and type-substitution stability give
\[
 R'[\alpha/a]\Rrightarrow_{\mathsf F}R^\star[\alpha/a]=M^\star.
\]
If the root is retained, inversion of the type-abstraction clause gives
\(N=(\Lambda a.R')[\alpha]\); induction followed by clause~7
contracts it to the same result.

The five outer constructors are exhaustive. In the root cases the premise
inductions allow any parallel reductions inside the body and, for term
application, the argument. Thus two term redexes, two type redexes, mixed
nestings, duplicated or erased arguments, disjoint positions and newly
exposed heads have all been included. Types themselves supply no additional
redex overlap. If \(M\Rrightarrow_{\mathsf F}N\) and
\(M\Rrightarrow_{\mathsf F}P\), the result just proved gives
\(N\Rrightarrow_{\mathsf F}M^\star\) and
\(P\Rrightarrow_{\mathsf F}M^\star\), proving the diamond property.
\end{proof}
\end{annexitem}

\begin{annexitem}{annex:p7.system-f.confluence}
 {thm:p07-system-f-confluence}{Finite joins for all raw combined reductions}
\begin{proof}
By the closure lemma, convert the two given finite ordinary reductions
from \(M\) into parallel paths, say of lengths \(m,n\). Construct
a rectangular array \(A_{i,j}\), \(0\leq i\leq m\),
\(0\leq j\leq n\). Let \(A_{0,0}=M\); put the first path along
the edge \(A_{i,0}\) and the second along \(A_{0,j}\).
For each \(i\geq1\), fill its row in increasing \(j\). The
already constructed terms \(A_{i,j-1}\) and \(A_{i-1,j}\) are
parallel successors of \(A_{i-1,j-1}\). The diamond property supplies
\(A_{i,j}\) as a parallel successor of both. The array is finite,
so this process finishes without any normalization hypothesis.

The lower and right edges give parallel paths from the two original
endpoints to \(A_{m,n}\). If either original path has length zero,
the other endpoint itself supplies the required join. Each finite parallel
path is an ordinary finite reduction by the other inclusion in the
closure lemma. This proves exactly the stated common-reduct property.

If two resulting terms are normal, neither has a genuine further
contraction. Their common reduct must therefore be each of them, modulo
the permitted alpha-renaming. This proves uniqueness of a normal form
when one exists, without proving its existence for arbitrary raw terms.
The separate strong-normalization theorem supplies existence for well-typed
inputs; it has not been used in this confluence proof.
\end{proof}
\end{annexitem}

\begin{annexitem}{annex:p7.system-f.type-syntax-decision}
 {thm:p07-system-f-type-syntax-decision}{Finite formation and bound-name comparison}
\begin{proof}
The formation procedure recurses only on proper type subexpressions.
At a variable its membership test is finite and agrees with the type-variable
rule. At an arrow the two recursive tests agree with the two formation
premises. At a universal type a fresh name can be found effectively by
searching the fixed name supply outside the finite set of names to avoid.
Consistent renaming changes no free variable. Checking the body in the
extended context is exactly the universal formation premise. Induction
therefore proves both termination and exactness. Equivalently, a parsed
type forms over \(\Delta\) exactly when all its free variables are
in that context.

The stack encoding of a type also traverses a finite tree. Renaming a
binder and all occurrences bound by it preserves every bound index and
every free-name token, so alpha-equal types have identical encodings.
Conversely, suppose the encodings agree. Traverse both types together,
and at each matching binder marker rename the two binders to the same
new name, fresh for all free names and for previously chosen binders.
At matching bound indices, the corresponding variable occurrences then
have the same chosen name. At free-name tokens they already have the
same free name. The arrow and universal constructor markers agree by
hypothesis. Structural induction shows that the two resulting named
types are literally identical, proving alpha equality of the originals.
In particular, a free variable is never confused with a bound index.

For substitution, a variable is either replaced by its specified finite
image or retained. An arrow substitutes in both children. At a universal
binder, its own bound occurrences are excluded; other substitutions proceed
under a name first freshened for their images. The recursion follows the
original type tree and does not revisit inserted images, so it terminates
even when an image itself contains a variable from the substitution domain.
The result is finite, since finitely many finite images are inserted.
The earlier type-formation substitution proof gives preservation of
formation under the stated source and target contexts.
\end{proof}
\end{annexitem}

\begin{annexitem}{annex:p7.system-f.synthesis-correct}
 {thm:p07-system-f-synthesis-correct}{Every synthesis case and every possible typing derivation}
\begin{proof}
Freshening is effective: the current finite contexts and finite syntax
contain only finitely many names. Alpha-renaming of a term binder changes
only its bound term occurrences. Alpha-renaming of a type binder changes
its bound type occurrences, including occurrences in annotations and
type arguments, and respects inner shadowing. Free names in the context
are not changed. Typing is preserved under these allowed bound renamings,
by the fresh-variable cases of the earlier typing-substitution lemma.
We may therefore prove the result for the consistently freshened input.

Every recursive synthesis call concerns a proper term subexpression.
The finite formation, alpha-equality and type-substitution procedures
terminate by the preceding lemma. New result types can be larger than
input annotations, but they are not passed back as new terms to synthesize.
Thus induction on term size proves termination of the whole procedure.
The recursive contexts stay valid: an annotation is checked before its
term variable is added, and a fresh type-variable extension preserves the
formation of all existing context types.

For soundness and completeness, strengthen the induction statement to
include well-formedness of a returned type and uniqueness up to alpha
equality among all types derivable for the same term and contexts.

At a variable, the unique context declaration is the only possible
variable-rule conclusion. A missing declaration cannot be supplied by
any other rule, since their conclusions have other term constructors.
Hence the returned type, or failure, is exact.

At \(\lambda x{:}\alpha.P\), successful formation of \(\alpha\)
and successful premise synthesis give a derivation for the body in
\(\Gamma,x:\alpha\). Arrow introduction gives exactly the returned
arrow type. Conversely, any derivation of this term must end in arrow
introduction with that same explicit annotation. Its premise exists
only if formation passes and body synthesis succeeds by induction.
The body's type is unique by induction, and therefore so is the arrow
result. Failure of either check precludes any such derivation.

At \(PN\), if synthesis returns an arrow for \(P\) and a matching
domain type for \(N\), induction supplies the two typing premises,
and arrow elimination gives the returned codomain. Conversely any
typing of \(PN\) must end in arrow elimination. By induction, the
type synthesized for \(P\) must be alpha-equal to that premise's
arrow type; its domain and codomain are consequently determined. The
type synthesized for \(N\) must match the domain. Thus every possible
derivation is accepted and has the same result. A universal type is
not an arrow type, and the rules contain no implicit instantiation that
could repair a failed shape check.

At \(\Lambda a.P\), after freshening, the sheet's two eigenvariable
conditions are the explicitly checked conditions. If body synthesis
under \(\Delta,a;\Gamma\) succeeds at \(\beta\), induction gives
its typing derivation and formation. Universal introduction gives
\(\forall a.\beta\), a type formed over \(\Delta\). Conversely,
the only possible last rule for this constructor is universal introduction,
with the same side conditions and a derivable premise. Body induction
therefore makes synthesis succeed, and every premise type is alpha-equal
to its result. Consistently closing the same binder preserves that equality.
This proves uniqueness and the rejection conclusion in this case.

At \(P[\alpha]\), successful formation of the explicit instance type
and synthesis of a universal function type provide the premises of
universal elimination. That rule's conclusion is exactly the computed
capture-avoiding instance. Formation is preserved by the type-substitution
lemma. Conversely, any typing derivation of this constructor must have
a universal function premise and that same supplied type argument.
Induction identifies the universal type up to alpha equality. Rename
its bound variable consistently before substitution; substituting in
alpha-equal types gives alpha-equal results by the same structural
substitution procedure. Thus the returned result is unique and agrees
with every derivation. A failed function-shape or formation check
precludes the required elimination rule.

All five constructors have now been treated. No reduction, normal form,
unification problem or guessing of missing type information occurs in
the argument.
\end{proof}
\end{annexitem}

\begin{annexitem}{annex:p7.system-f.church-checking}
 {thm:p07-system-f-church-checking}{Decision for the exact explicit judgment}
\begin{proof}
Take the finite contexts, explicit term and proposed type from the
statement. Validate their finite syntax and the proposed type's formation.
Run \(\mathsf{Synth}(\Delta,\Gamma,M)\). If it fails, its correctness
theorem rules out every typing derivation in those contexts, so reject.
If it returns \(\beta\), compare \(\beta\) and the proposed
\(\alpha\) by the terminating alpha-equality test. Accept exactly
when they agree. Soundness of synthesis and equality modulo the permitted
bound renamings prove that acceptance yields the required judgment.
Conversely, if that judgment is derivable, completeness and uniqueness of
synthesis force a successful return alpha-equal to \(\alpha\), so the
comparison accepts. Every part of this procedure terminates.

This decides only the supplied Church-style judgment, whose lambda
annotations, type abstractions and type arguments are explicit input.
It does not supply an algorithm to infer absent annotations or absent
instances for an erased Curry-style term. Neither confluence nor the
normalization theorem is a premise of this decision proof.
\end{proof}
\end{annexitem}
